\chapter{并排对照：同一根香肠的四种刀法}
\label{ch:side}

\section{总图}

\begin{figure}[htbp]
\centering
\begin{tikzpicture}[font=\footnotesize,
  row/.style={softbox=#1, minimum width=11.5cm, minimum height=1.05cm, align=left}]
  \node[row=Gen3Color] (h) at (0,3.6)
    {\textbf{Hom}：铃叮？ -- 珠串边变了？ -- 切};
  \node[row=Gen4Color] (c) at (0,2.3)
    {\textbf{Chain}：锁对齐？ -- 骨架变了？ -- 有环？ -- 切};
  \node[row=Gen5Color] (p) at (0,1.0)
    {\textbf{TopoPH}：铃叮？ -- 钉星 -- 翻面/连通滤网？ -- 切};
  \node[row=Gen51Color] (n) at (0,-0.3)
    {\textbf{Native}：炼 $Y$ -- 落网格钉星 -- 事件？ -- 切（无铃）};
\end{tikzpicture}
\caption{四种``发车信号''链路，一眼分清。}
\label{fig:four-pipelines}
\end{figure}

\section{谁在听指纹铃}

\begin{center}
\begin{tabular}{@{}lcccc@{}}
\toprule
 & Hom & Chain & TopoPH & Native \\
\midrule
Gear 指纹铃 & 是 & 否 & 是 & \textbf{否} \\
密码锁 & 否 & 是 & 否 & 否 \\
事件作真切点 & 否（第二闸） & 否（第二/三闸） & 否（滤网） & \textbf{是} \\
\bottomrule
\end{tabular}
\end{center}

\section{仓库隔离（再强调一遍）}

\begin{figure}[htbp]
\centering
\begin{tikzpicture}[font=\small]
  \foreach \i/\name/\col in {0/Hom仓/Gen3Color, 1/Chain仓/Gen4Color,
    2/PH仓/Gen5Color, 3/Native仓/Gen51Color} {
    \node[softbox=\col, minimum width=2.4cm] (r\i) at (\i*2.9,0) {\name};
  }
  \draw[<->, thick, dashed, gray] (r0)--(r1);
  \draw[<->, thick, dashed, gray] (r1)--(r2);
  \draw[<->, thick, dashed, gray] (r2)--(r3);
  \node[font=\footnotesize, gray, below=0.55cm of r1]
    {虚线 = 不能就地搬家；要换刀法就新开仓重切};
\end{tikzpicture}
\caption{四个房间，四套积木，别混着玩。}
\label{fig:repo-rooms}
\end{figure}

\section{速度直觉（非正式）}

\begin{metaphor}[厨房计时]
\begin{itemize}[nosep]
  \item \textbf{Chain}：多数字节只拧锁 —— 通常最快；
  \item \textbf{Hom}：多数字节只滚指纹 —— 接近 FastCDC；
  \item \textbf{TopoPH}：叮的时候还要算几何 —— 更贵；
  \item \textbf{Native}：几乎每字节炼 $Y$ —— 有结构性成本，靠 SIMD / 缓存标定打平体验。
\end{itemize}
具体数字看 bench，不看这本图画书。
\end{metaphor}
